\section{OBJETIVO}
Este arquivo tem por propósito ilustrar a estrutura padrão de saída de um documento técnico, exibindo a organização de seções, listas e itens de definição.

\noindent Não se trata de um inventário exaustivo de componentes, tampouco reproduz dados constantes em normas técnicas ou orientações detalhadas sobre processos produtivos e desenhos dimensionais.


\section{DOCUMENTOS DE REFERÊNCIA}
\begin{itemize}
\item \textbf{EMP-GER-ETG-001-0.1} – DOCUMENTO BASE A - ESPECIFICAÇÃO TÉCNICA GERAL DO EMPREENDIMENTO
\item \textbf{EMP-MEC-PRO-001-0.3} - DOCUMENTO BASE B - PROCEDIMENTOS DE SOLDAGEM
\item \textbf{EMP-GER-ETG-002-0.1} - DOCUMENTO BASE C - DIRETRIZES DE PROJETO DO CLIENTE
\end{itemize}

\section{ORIENTAÇÕES DE LEITURA}

\begin{enumerate}[label=\alph*)]
\item As fichas reúnem famílias de componentes segundo o material e o tipo cadastrado no banco de exemplo.

\item Cada linha da tabela identifica uma dimensão e apresenta os códigos de ambas as partes, a quantidade emitida e sua unidade.

\item Valores iguais a zero permanecem na planilha de entrada e são omitidos do documento. Quando uma família não tem figura cadastrada, a ficha usa a imagem padrão do template.
\end{enumerate}


\section{ESTRUTURA DE CODIFICAÇÃO}

\begin{figure}[H]
\centering

\begin{tikzpicture}[
    >=Stealth,
    font=\small
]

% Texto principal
\node[anchor=west] at (-5.35,0.1)
    {\large\underline{MEC}-\underline{AI}-\underline{10000}-\underline{000}};

% Linhas verticais
\draw (-4.65,-0.1) -- (-4.65,-2.6);
\draw (-3.75,-0.1) -- (-3.75,-1.9);
\draw (-2.75,-0.1) -- (-2.75,-1.2);
\draw (-1.55,-0.1) -- (-1.55,-0.5);

% Linhas horizontais + setas
\draw[->] (-1.55,-0.5) -- (0.5,-0.5);
\draw[->] (-2.75,-1.2) -- (0.5,-1.2);
\draw[->] (-3.75,-1.9) -- (0.5,-1.9);
\draw[->] (-4.65,-2.6) -- (0.5,-2.6);

% Textos
\node[anchor=west] at (0.6,-0.5)
    {Sequencial identificador único da dimensão característica};

\node[anchor=west] at (0.6,-1.2)
    {Sequencial identificador único da familia do elemento};

\node[anchor=west] at (0.6,-1.9)
    {Indicador de Material (Aço Inoxidável)};

\node[anchor=west] at (0.6,-2.6)
    {Indicador de área (Mecânica)};

\end{tikzpicture}

\end{figure}
\newpage
